% =============================================================================
% AGENTE  : Geneticista (Especialista em Domínio) — AUTOR PRINCIPAL DESTE MÓDULO
% MÓDULO  : Discussão — Português Brasileiro (discussion_pt.tex)
% NOTA    : Tradução de discussion_en.tex — o conteúdo técnico é de
%           responsabilidade do Geneticista; a tradução é do Cientista-Líder.
% STATUS  : RASCUNHO
% =============================================================================

\section{Discussão}
\label{sec:discussion}

Os resultados obtidos com o painel simulado de marcadores estão em boa
concordância com as expectativas estabelecidas na literatura de marcadores
moleculares. Os valores consistentemente mais elevados de \pic{} e \he{}
observados para os marcadores \ssr{} corroboram sua reconhecida superioridade
na captura de polimorfismo, característica atribuível ao modelo de mutação
em degraus que rege a evolução de repetições de microssatélites e ao
consequente acúmulo de múltiplos alelos por locus \autocite{Ellegren2004}.

Os marcadores \snp{}, embora individualmente menos informativos devido à sua
natureza bialélica, compensam por sua abundância no genoma. Plataformas de
genotipagem de alta densidade de \snp{} podem capturar padrões de
desequilíbrio de ligação em alta resolução, viabilizando estudos de
associação genômica ampla (GWAS) e seleção genômica com acurácia superior
à alcançável com painéis mais esparsos de \ssr{} \autocite{Goddard2007}. O
equilíbrio entre a informatividade por locus (\ssr{}) e a densidade de
marcadores (\snp{}) deve, portanto, ser avaliado conforme a aplicação
específica: análise de parentesco e identificação de cultivares tipicamente
favorecem o \ssr{}, enquanto predição genômica e mapeamento de QTLs em
escala favorecem matrizes densas de \snp{}.

A distribuição estreita dos valores de $H_o$ em torno de $H_e$ no conjunto
de dados simulado indica uma população em equilíbrio de Hardy--Weinberg
aproximado, conforme esperado dada a ausência de seleção, deriva ou
endogamia no modelo de simulação. Em populações reais, desvios desse
equilíbrio podem sinalizar fenômenos biológicos de interesse, como depressão
por endogamia, varreduras seletivas ou artefatos de genotipagem (por exemplo,
alelos nulos em locos \ssr{}), os quais devem ser sistematicamente testados
antes das análises subsequentes \autocite{Peakall2006}.

Em conjunto, esses resultados reforçam a complementaridade dos sistemas de
marcadores \snp{} e \ssr{}. O pipeline analítico modular desenvolvido aqui,
combinando computação estatística em Python com geração de documentos
\LaTeX{}, fornece um arcabouço reprodutível e extensível para estudos
futuros de caracterização de marcadores.

\subsection{Limitações e Perspectivas Futuras}

A presente análise é limitada pelo uso de dados simulados e por um modelo
bialélico para os locos \ssr{}, que subestima seu verdadeiro polimorfismo.
Trabalhos futuros devem incorporar dados de genotipagem empírica de
populações reais, empregar formulações multialélicas de \pic{} para
marcadores \ssr{} \autocite{Anderson1993}, e estender o arcabouço
estatístico para incluir estatísticas $F$ e análises de estrutura
populacional (\textit{e.g.}, STRUCTURE, ADMIXTURE).
